\documentclass[11pt,a4paper]{article}
\usepackage{fontspec}
\setmainfont{Noto Serif}
\setsansfont{Noto Sans}
\usepackage{babel}
\babelprovide[main,import]{russian}
\usepackage{amsmath,amssymb,mathtools,booktabs,geometry,fancyhdr,microtype,xcolor,hyperref}
\geometry{left=21mm,right=21mm,top=21mm,bottom=21mm}
\definecolor{ink}{HTML}{173D57}
\hypersetup{colorlinks=true,urlcolor=ink,linkcolor=ink,pdftitle={BGS-F. Затухание пеллева возврата и гейзенберговские времена}}
\pagestyle{fancy}\fancyhf{}\lhead{\small BGS-F}\rhead{\small Давление и интерференция траекторий}\cfoot{\thepage}
\setlength{\headheight}{15pt}\setlength{\parindent}{1em}\setlength{\parskip}{2.5pt}
\setlength{\emergencystretch}{3em}\tolerance=3000
\newcommand{\Tr}{\operatorname{Tr}}
\newcommand{\EE}{\mathbb E}
\newcommand{\R}{\mathbb R}
\newcommand{\Z}{\mathbb Z}
\newcommand{\C}{\mathbb C}
\newcommand{\si}{\mathrm{Sine}_2}
\newcommand{\eps}{\varepsilon}
\newcommand{\sinc}{\operatorname{sinc}}
\newcommand{\sect}[1]{\section*{\sffamily\color{ink}#1}}
\newcommand{\subt}[1]{\par\medskip\noindent\textbf{#1.}\quad}
\begin{document}
\noindent{\small\sffamily\color{ink} МАТЕМАТИЧЕСКАЯ ФИЗИКА\hfill 19 НОЯБРЯ 2026}
\begin{center}{\LARGE\sffamily Затухание пеллева возврата в BGS-F}\\[4pt]
{\large Давление неустойчивого якобиана и интерференция при $m\asymp N$}
\end{center}
\sect{Результат для бесконечной последовательности}
Для исходных матриц $U_N$, всех нечётных $N$ и фиксированного $\tau=10^{-80}$ положим $p_{m,N}=\Tr U_N^m$, $h=\log(3+2\sqrt2)$. Существуют числа $a>2/h$, $\eta>0$, $C<\infty$, зависящие от карты, такие, что
\[
 \boxed{\quad N^{-2}\sum_{1\le m\le a\log N}|p_{m,N}|^2\le C N^{-2\eta}.\quad} \tag{1}
\]
В частности, если $(1+\sqrt2)^k=h_k+v_k\sqrt2$, то
\[
 \frac{|p_{k,h_k}|}{h_k}=O(h_k^{-\eta}),\qquad k\to\infty\ \text{по нечётным }k. \tag{2}
\]
Предшествующий сертификат $|p_{209,h_{209}}|/h_{209}>1.8\cdot10^{-5}$ сохраняет силу. Теорема (2) отвечает на вопрос о его продолжении: на том же пеллевом времени $m=k$ не существует бесконечной подпоследовательности с нормированным следом, отделённым от нуля. Здесь важен строгий запас в $a>2/h$: граница $a<2/h$ не охватывает эти возвраты.

Для двух отмеченных общих участков траекторий при $N<m<2N$ получена также нижняя оценка знаковой суммы после удаления циклических пар. При суммировании с весом $w$ из (P) она равна снизу, в пределе, $\int_1^2 t(t-1)w(t)dt>0$. Кроме того, для явно выделенного множества нециклических пар без повторного счёта получена отрицательная асимптотика $-N(m/N)^3/6+o(N)$, а для класса одинаковых направленных рёбер доказан экспоненциальный рост. В этих теоремах используются оценки смешивания (7.7)--(7.8) исходной постановки. При $m=N=h_{209}$ для одной явной пары расположений общих участков получен отдельный строго положительный сертификат без этой зависимости.

\subt{Статус} Слабая сходимость к $\si$ и предел (P) остаются открытыми. Оценка (1) относится к логарифмическим временам. При гейзенберговских временах явно оценены положительные и отрицательные классы путей. Разделы 7--9 показывают, какие из них компенсируют друг друга и почему сумма по остальным парам остаётся необходимой.

\subt{Состав отчёта} Разделы 1--4 содержат доказательство (1), проверку нечётного квантования и вычисление вариации давления. Разделы 5--10 посвящены гейзенберговским временам и контрольному примеру. В приложении после этого текста воспроизведён полный отчёт с сертификатом для $h_{209}$. Нумерация формул и страниц приложения самостоятельна.
\newpage
\sect{1. Исходный оператор и его классическая карта}
Для нечётного $N$ обозначим $e_N(t)=e^{2\pi i t/N}$, чертой --- мультипликативное обращение по модулю $N$. В координате $x\in\Z/N\Z$
\[
 (T(a,b)u)(x)=e_N(b(x+\bar2 a))u(x+a),\quad
 M_{xy}=N^{-1/2}e_N(-\bar8(3x^2-2xy+3y^2)).
\]
Пусть $r(q)=(1-\cos14\pi q)(1+\cos2\pi q)$, $e^{i\alpha_0}=(-3+4i)/5$. Квантовый гамильтониан $A=\sum c_{ba}T(a,b)$ задан коэффициентами при $e^{2\pi i(bq+ap)}$ функции
\[
 f_0=\cos14\pi q+r(q)(1/10+\sin(2\pi q)/100+
                      \cos(2\pi p-\alpha_0)/100)-1/10 .
\]
В физическом базисе Вейля $W(a,b)=(-1)^{ab}T(a,b)$ из-за смешанного члена его символ имеет вид
\[
 g=\cos14\pi q+r(q)(1/10+\sin(2\pi q)/100)
      +r(q+1/2)\cos(2\pi p-\alpha_0)/100-1/10 .
\]
При $\hbar=(2\pi N)^{-1}$ матрица $U_N=e^{i\tau NA}M$ квантует
\[
 F_\tau=\phi^1_{-\tau g/(2\pi)}\circ B,\quad
 B=\begin{pmatrix}3&-4\\-2&3\end{pmatrix},\quad
 S=\begin{pmatrix}3&4\\2&3\end{pmatrix} .
\]
Здесь $\dot q=H_p$, $\dot p=-H_q$. Действительно, $M T(d)M^*=T(Sd)$ по прямому сравнению ядер. Для $W$ знаки совпадают, так как $S\equiv I\pmod2$. В физических координатах $W(a,b)$ даёт смещение $(-a,b)/N$. Следовательно, действующая на смещения матрица есть $B=\mathrm{diag}(-1,1)S\mathrm{diag}(-1,1)$. Неприводимость операторов сдвига и модуляции определяет метаплектический множитель с точностью до скаляра, который исчезает в модулях следов.

\subt{Проверка гиперболичности} Для всех $|u|\le10^{-80}$ дифференциал кика отличается от $I$ в евклидовой норме не более чем на $e^{10^4|u|}-1$. Достаточно просуммировать вторые производные рядов Фурье: $|a|\le1$, $|b|\le9$, сумма модулей коэффициентов $g$ меньше $1.2$, так что для гамильтониана $-g/(2\pi)$ верхняя граница есть $2\pi\cdot82\cdot1.2<10^4$. То же верно для обратного потока. В базисе из $(1,-1/\sqrt2),(1,1/\sqrt2)$ число обусловленности равно $\sqrt2$. Там $DF_u$ и $DF_u^{-1}$ отличаются от диагональных $DB$ и $DB^{-1}$ менее чем на $10^{-72}$ по норме. Конусы $|s|\le|u'|$, $|u'|\le|s|$ инвариантны строго вперёд и назад соответственно, с расширением доминирующей координаты больше 5. Конусный критерий даёт свойство Аносова на всём пути. По структурной устойчивости класс топологического сопряжения локально постоянен на этом связном пути. Поэтому $F_\tau$ сопряжена с $B$, топологически перемешивает и имеет энтропию $h$. Неустойчивое направление ориентируется конусом, его множители положительны.
\newpage
\sect{2. Строгое неравенство давления}
Пишем $v(x)=\log\|DF_\tau(x)|_{E^u(x)}\|$, $P(s)=P_{\rm top}(F_\tau,-sv)$. Потенциал гёльдеров. Замена нормы в одномерном неустойчивом расслоении меняет его на кограницу, то есть добавляет разность $b\circ F_\tau-b$. Из вариационного принципа, неравенства Рюэля и формулы энтропии Пезина для гладкой инвариантной площади имеем
\[
 P(0)=h,\qquad P(1)=0.
\]
При $s=1$ верхняя граница для всех инвариантных мер равна нулю, и мера площади её достигает. В частности, $P(1/2)>0$ по той же мере.

Две точки $z_0=(0,0)$ и $z_1=(0,1/2)$ неподвижны для $B$ и кика. В них $\nabla g=0$, $g_{pp}=g_{qp}=0$. Из $r''(0)=392\pi^2$, $r''(1/2)=8\pi^2$ и $\cos\alpha_0=-3/5$ следует
\[
 \frac{g_{qq}(z_0)}{\pi^2}=-\frac{19606}{125},\qquad
 \frac{g_{qq}(z_1)}{\pi^2}=-\frac{19594}{125}.
\]
В каждой точке линеаризация потока за единичное время есть точная матрица сдвига, поскольку квадрат его инфинитезимальной матрицы нулевой. Умножение на $B$ даёт
\[
 \Tr DF_\tau(z_i)=6-2\tau g_{qq}(z_i)/\pi .
\]
Определители равны единице. Следы различны и больше двух, поэтому различны и расширяющие собственные числа. Суммы $v$ вдоль этих двух орбит длины один различны. Тем самым $v$ не когомологичен константе.

Для перемешивающего диффеоморфизма Аносова давление семейства гёльдеровых потенциалов $-sv$ строго выпукло, когда $v$ не когомологичен константе [2,3]. Можно также рассуждать по случаю равенства: равенство выпуклости в середине вынудило бы равновесную меру среднего потенциала достигать обеих крайних вариационных границ. По свойству Гиббса для общей меры отношения весов шаров Боуэна на $n$ повторениях любой неподвижной точки ограничены сверху и снизу независимо от $n$. Тогда $n(v(z_i)-h)$ ограничено для обеих точек, что противоречит вычисленным множителям. Получаем строгое неравенство
\[
 0<p:=P(1/2)<h/2.                                                        \tag{3}
\]

Если $S_mv(x)=\sum_{j=0}^{m-1}v(F_\tau^jx)$, то для каждого $\epsilon>0$
\[
 \sum_{F_\tau^m x=x} e^{-S_mv(x)/2}\le C_\epsilon e^{(p+\epsilon)m}.      \tag{4}
\]
Действительно, множество неподвижных точек итерации есть разделённое по метрике Боуэна множество с одним и тем же параметром разделения: для периодических орбит близость в один период повторилась бы для всех целых времён, вопреки экспансивности. Определение давления через разделённые множества даёт (4). Этим объясняется выигрыш относительно простого подсчёта числа орбит.
\newpage
\sect{3. Следовая формула и нечётность $N$}
Применяем локальную логарифмическую следовую формулу Фора [1, теорема 12 в версии arXiv, (68)--(72)]. Для фиксированного $A_0>0$, любого $K>0$ и $1\le m\le A_0\log N$ она даёт сумму по $F_\tau^mx=x$ с фазами модуля один, знаменателем
\[
 2\sinh(S_mv(x)/2)
\]
и полуклассическим множителем $1+\sum_{s=1}^{J'}\hbar^s E_s(x,m)$, где $|E_s|\le C_s m^s$ для фиксированного порядка; суммарный остаток $O_{A_0,K}(N^{-K})$. Порядок выбирается после $A_0,K$. Число $e^{S_mv(x)}$ есть неустойчивый мультипликатор на периоде. Гиперболичность даёт $S_mv(x)\ge c m-C$; для малых $m$ используем конечность числа орбит и строгую гиперболичность. Поэтому каждый модуль амплитуды ограничен $C e^{-S_mv(x)/2}$.

\subt{Решёточный знак} У Фора в формулах (21)--(22), (35) выбрано чётное $N$ для устранения знака целочисленных трансляций. Приведём переход для матриц настоящей задачи. Пусть $D_n$ --- плоское вейлево смещение на $n\in\Z^2$. Тогда $D_nD_l=(-1)^{N(n_1l_2-n_2l_1)}D_{n+l}$. Положим $L_n=(-1)^{N n_1n_2}D_n$. Имеем
\[
 L_nL_l=L_{n+l},\qquad \widehat U L_n\widehat U^{-1}=L_{Bn}.
\]
Второе равенство следует из точной метаплектической ковариантности, $B\equiv I\pmod2$ и периодичности гамильтониана кика. Сектор $L_{e_j}=1$ состоит из периодических гребёнок на $x/N$ и реализует данные $W(a,b)$. Проектор периодизации в смысле распределений есть $\sum_n L_n$, как произведение сумм по двум генераторам. В доказательстве торической следовой формулы предварительные слагаемые имеют вид $L_{-n}\widehat U^m$. Сопряжение посредством $L_l$ переводит их в $L_{-(n+(B^m-I)l)}\widehat U^m$ \emph{точно}. Тем самым группировка по $\Z^2/(B^m-I)\Z^2$ и разворачивание интегралов по фундаментальной ячейке остаются в силе. След плоского слагаемого для представителя отличается от локального слагаемого с $D_{-n}$ одним постоянным знаком. Оценки микролокализации, нормальной формы и её остатка в [1, разделы 4--6] для плоского слагаемого сохраняют свои модули. Это доказывает применённую выше оценку и для нечётного $N$. Дополнительный скаляр метаплектического выбора также имеет модуль один.

Из (4) теперь следует $|p_{m,N}|\le C_\epsilon e^{(p+\epsilon)m}+O(N^{-K})$. Выберем $0<\epsilon<h/2-p$, $2/h<a<1/(p+\epsilon)$, $\eta=1-a(p+\epsilon)>0$. Суммирование геометрической прогрессии даёт (1), с константой для выбранных параметров. При $N=h_k$ имеем $k=\frac2h\log(N+\sqrt{N^2+1})$, следовательно $k\le a\log N$ начиная с некоторого $k$. Это доказывает (2) для исходного фиксированного $\tau$. Вклад этого временного диапазона в (P) стремится к нулю: вес ограничен, а вклад вычитаемой рампы равен $O((\log N)^2/N^2)$.
\newpage
\sect{4. Величина выигрыша в первом ненулевом порядке}
Обозначим $p(u)=P_{\rm top}(F_u,-v_u/2)$ при переменном $u$ в окрестности нуля. Для этого вспомогательного семейства
\[
 h/2-p(u)=\frac{7834656061}{32000000}\pi^2u^2+o(u^2).                   \tag{5}
\]
Строгое неравенство (3) при $u=10^{-80}$ было получено непосредственно из двух неподвижных точек. В (5) остаток не снабжён численным порогом.

Перенесём потенциал на фиксированную карту $B$ топологическим сопряжением. Нормируем неустойчивый вектор $r_u$ условием $\ell r_u=1$, где $r=(1,-1/\sqrt2)^{\mathsf T}$, $\ell=(1/2,-1/\sqrt2)$, $Br=e^hr$, $\ell B=e^h\ell$. В этой нормировке логарифм расширения в перенесённых координатах равен $h+u b+O(u^2)$ в некотором $C^\beta$, где
\[
 b(x)= -\frac{g_{qq}-g_{pp}/2}{2\pi\sqrt2}(Bx).
\]
Действительно, дифференциал кика в первом порядке равен $I-uJ\,\mathrm{Hess}(g)/(2\pi)$, $J=\left(\begin{smallmatrix}0&1\\-1&0\end{smallmatrix}\right)$. В производной множителя $\ell DF_u r_u$ вклад изменения $r_u$ исчезает вследствие нормировки. При $u=0$ базовый дифференциал постоянен, поэтому сдвиг базовой точки также исчезает.

Для обоснования нормы остатка выберем малое $\beta>0$ с $\max(\mathrm{Lip}(B),\mathrm{Lip}(B^{-1}))^\beta<e^h$. Запишем сопряжение как $x\mapsto x+w(x)$. Линеаризация уравнения сопряжения по $w$ есть $w\circ B-Bw$. На стабильной компоненте её обратный получается суммированием назад по $B$, на нестабильной --- вперёд; члены оцениваются геометрической прогрессией со знаменателем меньше единицы в $C^\beta$. Нелинейная часть есть подстановка $x+w(x)$ в гладкую периодическую функцию, дважды дифференцируемая на этом пространстве. Теорема о неявной функции и единственность близкого сопряжения дают $w=O(u)$ в $C^\beta$. Для наклона неустойчивого вектора в этих координатах линеаризация граф-преобразования равна $e^{-2h}$, умноженному на оператор композиции с $B^{-1}$; тем же аргументом $r_u=r+O(u)$ в $C^\beta$. Этих оценок достаточно для указанного разложения.

Для давления на фиксированной карте первая вариация есть интеграл по равновесной мере, вторая при нулевом потенциале --- асимптотическая дисперсия по мере максимальной энтропии, здесь по площади [2,4]. Положим $z=v_u\circ H_u-h$ в принятой нормировке, $\mathcal P$ --- давление на $B$. Разность $[\mathcal P(0)+\mathcal P(-z)]/2-\mathcal P(-z/2)$ равна $h/2-p(u)$, так как $P_{F_u}(-v_u)=0$. Линейные члены сокращаются, и остаётся $u^2\sigma^2(b)/8+o(u^2)$.

Уберём общий сдвиг аргумента $B$ в $b$. Его коэффициенты, делённые на $\sqrt2\pi$, равны $(j^2-l^2/2)\widehat g(j,l)$. При корреляции на шаг $n$ частота переходит в $(B^\mathsf T)^n(j,l)$. Для $|n|\ge2$ общих частот нет: модуль второй координаты при $j\ne0$ не меньше $2v_{2|n|}-h_{2|n|}\ge7$, а при $j=0,l\ne0$ он не меньше $h_{2|n|}$. Исходно $|l|\le1$. Точное суммирование даёт ковариации на шагах 0 и 1, делённые на $2\pi^2$, равные $313386243/320000$ и $-7/8000000$. Учитывая оба знака шага 1, получаем (5).
\newpage
\sect{5. Следы на движущихся окнах и общие участки}
Для $T=\lfloor a\log N\rfloor$ из геометрической оценки раздела 3 следует $\sum_{d=1}^T|p_{d,N}|/N=O(N^{-\eta})$. В пространстве Гильберта--Шмидта возьмём $T$ единичных векторов $U^{r+j}/\sqrt N$, $1\le j\le T$, где $r$ произвольно. Их матрица Грама имеет диагональ один и сумму модулей внедиагональных элементов каждой строки $O(N^{-\eta})$. Её норма не больше $1+O(N^{-\eta})$. Применяя соответствующее рамочное неравенство к $I/\sqrt N$, получаем
\[
 \sup_{r\in\Z}\sum_{m=r+1}^{r+T}\frac{|p_{m,N}|^2}{N^2}\le1+O(N^{-\eta}).\tag{6}
\]
Поэтому для фиксированного $b>0$ среди $N\le m\le2N$ имеется $O(N/(b^2\log N))$ времён с $|p_{m,N}|\ge bN$. Даже один такой пик может быть заметен в (P), поэтому исключение большинства времён оставляет существенную возможность.

\subt{Две траектории} Для замкнутого слова $x=(x_0,\ldots,x_{m-1})$ обозначим $A(x)=\prod_j U(x_j,x_{j+1})$, индексы читаются циклически. Вес пары $(x,y)$ есть $A(x)\overline{A(y)}$. Ниже каждый путь имеет отмеченное начало, как в разложении $|\Tr U^m|^2$. Положим $P_{ab}=|U_{ab}|^2$, $J_{ab}=1/N$, $\rho=\|P-J\|_{2\to2}$, $\beta=\max P_{ab}$. Для бесконечной последовательности в разделах 5--7 используем как вход утверждения (7.7)--(7.8) приложенной постановки: для всех достаточно больших нечётных $N$ имеем $\rho\le1/4$, $\beta\le K_0/N$ с постоянной $K_0$. Для конечного сертификата эти оценки будут непосредственно проверены.

Зафиксируем $t$ общих рёбер в начале обоих слов, то есть $x_j=y_j$ для $0\le j\le t$. Второй общий участок имеет те же $t$ рёбер, но начинается в $x$ после промежутка длины $g$, в $y$ после промежутка $g'$. Пишем $D=m-2t+1$, $0\le g,g'<D$, $h_1=D-1-g$, $h'_1=D-1-g'$. Промежуток длины ноль допускается. Пусть $C_{g,g'}$ --- полная сумма по таким парам, $C^{\rm nc}_{g,g'}$ --- сумма после исключения всех циклических вращений. Введём матрицу $K_{u,v}(i,j)=(U^u)_{ij}\overline{(U^v)_{ij}}$. Суммирование по внутренним вершинам даёт точное равенство
\[
 C_{g,g'}=\Tr(P^t K_{g,g'} P^t K_{h_1,h'_1}).                          \tag{7}
\]
Действительно, фиксируем четыре конца общих участков. На совпадающем участке пары амплитуд перемножаются в цепь $P$; в каждом промежутке две суммы независимы при данных концах, что даёт $K_{u,v}$. Осталось замкнуть четыре индекса. Каждая строка и столбец матрицы $|K_{u,v}|$ имеют сумму не более единицы по Коши и унитарности. Поэтому $\|K_{u,v}\|\le1$.
\newpage
\sect{6. Вычисление знаковой суммы при $m\asymp N$}
Запишем $P^t=J+E$, $\|E\|\le\rho^t$. При двух множителях $J$ выражение (7) равно $|p_{g-g',N}|^2/N^2$, поскольку $J$ --- одномерный проектор, а сумма элементов $K_{u,v}$ есть $\Tr U^{u-v}$. Два слагаемых с одним $E$ ограничены каждое числом $\rho^t$, слагаемое с двумя $E$ --- числом $N\rho^{2t}$. Итак,
\[
 \left|C_{g,g'}-|p_{g-g',N}|^2/N^2\right|\le2\rho^t+N\rho^{2t}.         \tag{8}
\]

При $g\ne g'$ циклическая пара, удовлетворяющая ограничениям, вынуждает повторение одного и того же блока из $t+1$ вершин в $x$ на двух разных позициях. Иначе первая отметка определяет нулевой сдвиг, а вторая требует ненулевого. Для фиксированных позиций ориентируем круг так, чтобы расстояние вперёд между началами было $1\le d\le m/2$. При $t<m/3$ обе копии помещаются в разрезанное слово. Сумма положительных марковских весов $|A(x)|^2$ для этих совпадений не превосходит $N\beta^{t+1}$: старт выбирается $N$ способами, на каждом из $t+1$ шагов второй копии предписана уже посещённая вершина, остальные шаги суммируются стохастически. Условие замыкания можно отбросить сверху. Этот довод охватывает и перекрывающиеся копии. Не более $m^2$ пар начал и $m$ вращений дают
\[
 0\le C_{g,g'}-C^{\rm nc}_{g,g'}\le N m^3\beta^{t+1}\quad(g\ne g').      \tag{9}
\]
Рассмотрим $N\le m\le2N$, $t=\lceil12\log N\rceil$. Сумма правых частей (8),(9) по всем $g\ne g'$ равномерно равна $o(1)$, даже до нормировки. Поэтому для $\mathcal B_{m,N}=\sum_{g\ne g'} C^{\rm nc}_{g,g'}$ получаем
\[
 \mathcal B_{m,N}
 =\frac{D^2}{N}-D+
   \frac1{N^2}\sum_{j\ne l}\left|\sum_{d=0}^{D-1}
                  e^{2\pi i d(\theta_j-\theta_l)}\right|^2+o(1).          \tag{10}
\]
Здесь мы прибавили в сумму по $g,g'$ ровно $D$ диагональных вкладов единичной величины, затем раскрыли квадрат следа и выделили $j=l$. Функция $\mathcal B$ вещественна благодаря перестановке слов и $g,g'$. Последняя явная сумма в (10) неотрицательна.

Умножим на гладкий вес $w(t')=|\widehat b_*(t')|^2$ из (P). Переход к интегралу в первых двух слагаемых равномерен, так как $D=m+O(\log N)$. Доказано
\[
 \boxed{\liminf_{N\to\infty\ {\rm odd}}\frac1{N^2}
 \sum_{N<m<2N}w(m/N)\mathcal B_{m,N}
 \ \ge\ \int_1^2 t'(t'-1)w(t')\,dt'\ >\delta_* .}                     \tag{11}
\]
Здесь $\delta_*=\int_1^2(t'-1)w(t')dt'>0$. Последнее строгое сравнение следует из $w>0$ в $(1,2)$. В (11) знаки и фазы комплексных амплитуд удержаны до суммирования.
\newpage
\sect{7. Компенсация и предел применимости}
Пусть $A_0$ --- условие совпадения начального участка длины $t$, $M_2(x,y)$ --- число способов отметить второй участок с $g\ne g'$. Тогда $\mathcal B$ есть сумма весов нециклических пар с множителем $\mathbf1_{A_0}M_2$. Обозначим через $F_{m,N}$ аналогичную сумму только с $\mathbf1_{A_0}$. До удаления вращений она равна $\Tr(P^t K_{m-t,m-t})=1+O(N\rho^t)$. Совпадающие слова дают $\Tr P^m=1+O(N\rho^m)$, прочие вращения оцениваются по повторению блока, как в (9). Поэтому $F_{m,N}=o(1)$ равномерно на полосе. Для остатка с весом $\mathbf1_{A_0}(1-M_2)$ имеем точное равенство
\[
 R_{m,N}=F_{m,N}-\mathcal B_{m,N}.                                       \tag{12}
\]
Таким образом, (11) вынуждает отрицательную компенсацию величиной по крайней мере правой части (11) после того же суммирования. Вес $1-M_2$ меняет знак и учитывает повторные отметки. В исходном $Y_{m,N}$ все нециклические пары имеют кратность один. Формула (12) контролирует компенсацию внутри класса с заданным начальным совпадением; вклад пар вне этого класса остаётся неизвестным. В частности, (12) не устанавливает требуемую константу $-\delta_*$ в (P$'$).

\subt{Любой один синхронный участок} На фиксированном взаимном сдвиге можно снять выбор начальной позиции. Для каждого начала $j$ отметим совпадение $t$ рёбер. Для непустого набора обязательных начал $S$ обозначим через $L(S)$ число рёбер в объединении блоков. Сумма весов по парам с этими совпадениями есть след произведения $P^{\ell}$ на общих компонентах и $K_{u,u}$ в промежутках. Все множители дважды стохастичны и сохраняют константы. В их ортогональном дополнении норма произведения ограничена $\rho^{L(S)}$. Поэтому любая такая совместная сумма отличается от единицы не более чем на $N\rho^{L(S)}$.

По обращению Мёбиуса суммарная ошибка по модулям на всех \emph{непустых точных наборах} отмеченных начал, по отношению к единичной массе на наборе всех начал, ограничена $N\sum_{S\ne\varnothing}2^{|S|}\rho^{L(S)}$. Если блоки покрывают весь круг, их вклад в последнюю сумму без $N$ не больше $(3\rho)^m$. Иначе каждая компонента длины $l\ge t$ с заданным левым концом допускает суммарный вес $2^{|S|}$ не более $3^{l-t+1}$. Выбор $r$ левых концов стоит не более $m^r/r!$. Суммирование длин даёт общую оценку
\[
 N\left[(3\rho)^m+
 \exp\{3m\rho^t/(1-3\rho)\}-1\right]=o(1).
\]
В полном наборе начал слова совпадают. Из непустых собственных наборов можно удалить их остальные циклические пары с суммарной ошибкой из (9). Отсюда нециклическая знаковая сумма по всем парам, имеющим хотя бы один общий длинный блок \emph{в заданном выравнивании}, есть $o(1)$. При объединении разных выравниваний появляются как раз многократно отмеченные пересечения типа (10). Контроль всех этих пересечений на полной полосе не получен.
\newpage
\sect{8. Отрицательная сумма без кратности отметок}
Далее строится \emph{подмножество самих пар} в $Y_{m,N}$. Для $N\le m\le2N$, $c=m/N$, $t=\lceil12\log N\rceil$ сохраняем начальное совпадение $A_0$ из раздела 7. Пусть $Q(x)=|A(x)|^2$. Его суммарная масса $Z=\Tr P^m=1+O(N\rho^m)$. Для каждой тройки $i_1<i_2<i_3$ с $x_{i_1}=x_{i_2}=x_{i_3}=v$ разрежем круг в этих трёх местах. Получаются три пути из $v$ в $v$. Поменяем циклический порядок трёх путей на другой. Пару включаем в множество $E_{m,N}$, если выполнены два условия.
\begin{itemize}
\item Все три разреза лежат вне $[0,t]$. Полученное слово ставим началом в скопированное место исходного индекса 0. Так сохраняется $A_0$.
\item Каждое из трёх выделенных входящих рёбер $(x_{i_j-1},v)$ и каждое выделенное исходящее ребро $(v,x_{i_j+1})$ встречается в исходном слове ровно один раз.
\end{itemize}
Все направленные рёбра сохранены с кратностями, так что $A(y)=A(x)$. В вершине $v$ исчезают ровно три старые тройки «предшественник, $v$, преемник». Три новые переходят к другим преемникам. Выделенные входящие и исходящие рёбра попарно различны в своих списках. Поэтому разность двух мультимножеств троек восстанавливает три входящих ребра, а их уникальность в $x$ --- исходные позиции. В остальных вершинах переходы неизменны. Это доказывает одновременно нецикличность пары и отсутствие повторного счёта, в том числе при дополнительных посещениях $v$.

\subt{Масса троек} Для фиксированной вершины $v$ и $r$ различных обязательных моментов с круговыми промежутками $d_j\ge1$ точная сумма $Q(x)$ равна $\prod_{j=1}^r(P^{d_j})_{vv}$. При всех $d_j\ge t$ это $N^{-r}(1+O_r(N\rho^t))$. Вообще $(P^d)_{ab}\le K_0/N$ при $d\ge1$, по стохастичности и оценке для $P$. Наборов с промежутком меньше $t+3$ или с разрезом в $[0,t]$ всего $O_r(m^{r-1}(t+3))$. Их суммарная масса для данного $v$ есть $O_r(t/N)$. До требования уникальности это даёт массу троек $\binom m3/N^3+O(t/N+N\rho^t)$ на вершину.

Если одно из выделенных рёбер встретилось ещё раз, то совпали входящие или исходящие соседи в двух обязательных посещениях $v$. Учитываем $r=3$, если второе посещение уже выделено, и $r=4$, если оно дополнительное. Близкие посещения уже оценены. Для далёких посещений фиксируем соседей $a_j,b_j$ до и после $v$. Сумма имеет вес $\prod_j P_{a_jv}P_{vb_j}(P^{d_j-2})_{b_j a_{j+1}}$. Длинные переходы ограничены $(2/N)^r$ в произведении. При совпадении, например, $a_j=a_l$ сумма по соседям оставляет множитель $\sum_aP_{av}^2\le K_0/N$; остальные суммы не больше единицы. После подсчёта моментов вклад $O_r(1/N)$ на вершину. Конечное число вариантов выбора совпавших меток включено в константу.

\subt{Целочисленная проверка разбиения} Для матрицы Фурье размера 3 веса пар лежат в $3^{-m}\mathbb Z[e^{2\pi i/3}]$. Прямой перебор всех слов при $t=1$ даёт следующую проверку определений. «Остальные» обозначает все нециклические пары с заданным начальным ребром за вычетом $E_{m,3}$.
\begin{center}\small\begin{tabular}{@{}rrrrr@{}}\toprule
$m$ & Все с общим началом & Циклические & $E_{m,3}$ & Остальные \\ \midrule
7 & 1 & $403/243$ & $8/243$ & $-56/81$\\
9 & 1 & $4123/2187$ & $320/6561$ & $-6128/6561$\\ \bottomrule
\end{tabular}\end{center}
Здесь используется группировка фаз $1,\omega,\omega^2$ и равенство $|a+b\omega+d\omega^2|^2=a^2+b^2+d^2-ab-ad-bd$. Вторая проверка подтверждает восстановление исходных разрезов перебором 6561 трёхбуквенного слова длины 8. Числа служат контролем подсчёта и не участвуют в предельном доказательстве.

\newpage
\sect{9. Знак, размер и экспоненциальная компенсация}
Суммируя по всем вершинам, получаем равномерно для $N\le m\le2N$
\[
 \sum_{(x,y)\in E_{m,N}} A(x)\overline{A(y)}
   =N c^3/6+o(N).
\]
Вычтем это \emph{множество} из всех нециклических пар с общим началом $A_0$. Для оставшегося множества $E^-_{m,N}$ из оценки $F_{m,N}=o(1)$ раздела 7 следует
\[
 \boxed{\sum_{(x,y)\in E^-_{m,N}} A(x)\overline{A(y)}
      =-N c^3/6+o(N).}                                                     \tag{13}
\]
Следовательно, его взвешенный предел в нормировке (P$'$) на $N<m<2N$ равен $-\int_1^2 c^3w(c)dc/6<-\delta_*$. Сравнение строго, поскольку минимум $c^3/6-c+1$ на $[1,2]$ равен $1-2\sqrt2/3>0$. Это действительный отрицательный вклад с единичными коэффициентами каждой пары. Сумма $E_{m,N}\cup E^-_{m,N}$ равна $o(1)$ на время: их компенсация не вычисляет $Y_{m,N}$ вне $A_0$.

\subt{Все пары с одинаковыми направленными рёбрами} Пусть $H_{m,N}$ --- сумма по всем нециклическим парам с одинаковыми мультимножествами направленных рёбер. Каждое слагаемое равно $Q(x)\ge0$. Для любой последовательности $m/N\to c\in[1,2]$ имеет место
\[
 \liminf \frac1N\log H_{m,N}\ \ge \sum_{r=3}^\infty e^{-c}\frac{c^r}{r!}\log((r-1)!)>0.                 \tag{14}
\]
Для доказательства под вероятностью $Q/Z$ обозначим число посещений $v$ через $n_v$. Его биномиальный момент порядка $r$ по формуле раздела 8 стремится к $c^r/r!$, равномерно по $v$ для каждого фиксированного $r$. Все моменты сверху не больше $(2K_0)^r/(r! Z)$ для $r\ge1$. Ряд $\EE z^{n_v}=\sum_r\EE\binom{n_v}{r}(z-1)^r$ сходится равномерно на компактах, и его предел есть $e^{c(z-1)}$. Сходимость коэффициентов у нуля даёт пуассоновский закон. Для фиксированного $r$ вероятность $n_v=r$ при наличии повторяющегося входящего или исходящего соседа есть $O_r(t/N)$ по уже доказанной оценке с $r$ разрезами.

Зафиксируем $L\ge3$. В каждой вершине с $3\le n_v=r\le L$ и различными входящими, а также различными исходящими рёбрами, можно заменить круговой порядок $r$ экскурсий любым из $(r-1)!$ порядков. Операции выполняем последовательно по номерам вершин. Одна операция сохраняет связность обхода, мультимножество рёбер и все переходы в других вершинах. Два разных выбора в текущей вершине задают разные соответствия «вход--выход», различимые и после всех последующих операций. По первому различающемуся выбору получаем инъекцию в циклические классы слов. Тождественные выборы дают один исходный класс. Итак, $H_{m,N}\ge Z(\EE e^{S(x)}-1)$, где $S$ есть сумма $\log((r-1)!)$ по этим вершинам. По Йенсену и равномерному пуассоновскому пределу $N^{-1}\log H$ ограничен в пределе снизу частичной суммой (14). Увеличим $L$.

Полная сумма по остальным нециклическим парам равна $-H_{m,N}+O(N^2)$, поскольку $|p_{m,N}|^2\le N^2$ и циклическая часть не больше $mN$. Так установлена экспоненциально точная \emph{относительная} компенсация двух классов на всей полосе. Для (P) нужна абсолютная точность после суммирования на существенно меньшем масштабе. Экспоненциальное разрастание классов одинаковых обходов известно и в задачах о квантовых графах [5]; в (14) указан показатель снизу для данного режима роста размерности.
\newpage
\sect{10. Конечная проверка и контрольный пример}
При $N=h_{209}$ оценку смешивания можно доказать напрямую. У матрицы $A$ вне диагонали только два элемента на строку с модулем $|r|/200\le 1/50$. Отделив диагональ $A_0$, получаем в представлении взаимодействия $\|e^{i\xi A}-e^{i\xi A_0}\|_{\infty\to\infty}\le e^{0.04|\xi|}-1=:u$, где $\xi=\tau N<0.501$. Это следует из ряда Дайсона, поскольку диагональные фазы не меняют строчные суммы модулей. Так как $|M_{xy}|=N^{-1/2}$, получаем $|P_{xy}-1/N|\le(2u+u^2)/N$, $\rho<1/20$, $\beta<2/N$. Возьмём $m=N$, $t=100$, $g=\lfloor(m-2t)/2\rfloor$, $g'=g-209$. Тогда (8),(9) и сертификат приложения дают
\[
 \Re C^{\rm nc}_{g,g'} >3.23\cdot10^{-10}.
\]
Это нижняя оценка суммы амплитуд для конкретного расположения отметок непосредственно при $m=N$; нормировка здесь соответствует (7), без деления на $N$. Неравенство не постулирует поведение других расположений.

\subt{Проверка достаточности полученных оценок} На простых размерах $p>3$ обозначим $X u(x)=u(x+1)$, $D_{xx}=e_p(x^3)$, а через $\mathcal F$ --- унитарное дискретное преобразование Фурье. Положим $V_p=\mathcal F D X D^*\mathcal F^*$. Его спектр состоит из всех $p$-х корней единицы. Матрица $DXD^*$ есть $X$ с квадратичной фазой $x^3-(x+1)^3$. После преобразования Фурье каждый элемент имеет модуль $1/\sqrt p$: квадрат модуля соответствующей квадратичной суммы Гаусса равен $p$, поскольку суммирование по общей переменной зануляет все ненулевые сдвиги разности. Здесь старший коэффициент $-3$ обратим. Следовательно, у примера $P=J$ и $\Tr V_p^m=0$ для $1\le m<p$. При $m=p$ имеем $\Tr V_p^p=p$. Выполняются и коротковременная оценка (1), и сравнения типа (8)--(11).

Для исходного веса (P) предел теста в этом примере равен $w(1)-\int_0^2\min(t,1)w(t)dt<-10^{-3}$. Проверка знака воспроизводится рационально. При $\phi(t)=\exp(-1/(1-t^2))\mathbf1_{|t|<1}$ имеем $w=(\phi*\phi)^2/(4(\int\phi)^4)$, и $w$ убывает на положительной полуоси по представлению через наложение симметричных интервалов. Поэтому достаточно вычесть $w(1/2)/8+5w(3/4)/32$ из $w(1)$. Разбив $[-1,1]$ на 256 ячеек, для $e^{-x}$ на узлах используем $(1-x/512)^{512}\le e^{-x}\le(1+x/512)^{-512}$ при $0\le x<512$, с нулевым значением $\phi$ на границе поддержки. Округление наружу на сетку $2^{-36}$ и суммы Дарбу дают для этой верхней границы число меньше $-0.00294$. Таким образом, полученных ограничений на спектр и на модули элементов недостаточно для вывода (P). Контрольный оператор введён только для проверки логического перехода. Для $U_N$ остаётся исследовать неотмеченные фазы при $m\asymp N$.

\subt{Источники} [1] F. Faure, \href{https://doi.org/10.5802/aif.2341}{Ann. Inst. Fourier 57 (2007), 2525--2599}, \href{https://arxiv.org/abs/nlin/0610004}{arXiv:nlin/0610004v2}. [2] R. Bowen, \href{https://doi.org/10.1007/978-3-540-77695-6}{Equilibrium States and the Ergodic Theory of Anosov Diffeomorphisms}. [3] \href{https://tamara.ccny.cuny.edu/publications/AnosovDiffeos.pdf}{Realization of Anosov diffeomorphisms on the torus}, §1, о строгой выпуклости. [4] W. Parry, M. Pollicott, \emph{Zeta functions and the periodic orbit structure of hyperbolic dynamics}, Astérisque 187--188, гл.~4. [5] G. Tanner, \href{https://doi.org/10.1088/0305-4470/33/18/304}{Spectral statistics for unitary transfer matrices of binary graphs} (2000). В архиве приложены исходная постановка, предыдущий сертификат, исходники и проверки обоих исследований.
\end{document}
